\subsection{Bilinear Form Types and Congruence Invariants}\label{subsec:metric-vector-spaces-and-bilinear-forms}

% **Role:** arbitrary-field metric geometry.
%
% Symmetric, alternating, and skew-symmetric bilinear forms; rank, discriminant,
% and canonical forms under congruence.

Metric vector spaces retain the algebraic structure of inner-product geometry without imposing positivity.\\
Unless stated otherwise, $\mathbb{F}$ is an arbitrary field and every vector space is finite-dimensional.\\
General bilinear pairings, their radicals, their matrices, and their basis-change law were established earlier.\\
Here the two input spaces coincide, and a fixed bilinear form $B$ is written $B(\mathbf{x},\mathbf{y})=\innerprod{\mathbf{x},\mathbf{y}}$.

\begin{definition} \label{def:metric-vector-space-basic-form-types} \hlt{Symmetric, Skew-Symmetric, and Alternating Forms}\\
Let $\mathcal{V}$ be a vector space over $\mathbb{F}$ and let $B:\mathcal{V}\times\mathcal{V}\to\mathbb{F}$ be bilinear.
\begin{enumerate}[label=(\roman*)]
\setlength{\itemsep}{0pt}
\item $B$ is a \hlt{bilinear form} on $\mathcal{V}$; when fixed, write $B(\mathbf{x},\mathbf{y})=\innerprod{\mathbf{x},\mathbf{y}}$.
\item $\operatorname{Sym}(B)$ means $(\forall\,\mathbf{x},\mathbf{y}\in\mathcal{V})\,B(\mathbf{x},\mathbf{y})=B(\mathbf{y},\mathbf{x})$; then $B$ is \hlt{symmetric}.
\item $\operatorname{Skew}(B)$ means $(\forall\,\mathbf{x},\mathbf{y}\in\mathcal{V})\,B(\mathbf{x},\mathbf{y})=-B(\mathbf{y},\mathbf{x})$; then $B$ is \hlt{skew-symmetric}.
\item $\operatorname{Alt}(B)$ means $(\forall\,\mathbf{x}\in\mathcal{V})\,B(\mathbf{x},\mathbf{x})=0$; then $B$ is \hlt{alternating}.
\item A \hlt{metric vector space} is a pair $(\mathcal{V},B)$ such that $\operatorname{Sym}(B)$, $\operatorname{Skew}(B)$, or $\operatorname{Alt}(B)$ holds.\\
The symmetric and alternating cases define \hlt{orthogonal geometry} and \hlt{symplectic geometry}, respectively.
\end{enumerate}
\end{definition}

\begin{proposition} \label{prop:form-type-relations-by-characteristic} \hlt{Characteristic Controls the Form Types}\\
Let $B(\mathbf{x},\mathbf{y})=\innerprod{\mathbf{x},\mathbf{y}}$ be a bilinear form on $\mathcal{V}$.
\begin{enumerate}[label=(\roman*)]
\setlength{\itemsep}{0pt}
\item $\operatorname{Alt}(B)\Rightarrow\operatorname{Skew}(B)$.
\item If $\operatorname{char}(\mathbb{F})\neq2$, then $\operatorname{Sym}(B)\wedge\operatorname{Alt}(B)\Rightarrow B=0$.
\item If $\operatorname{char}(\mathbb{F})\neq2$, then $\operatorname{Alt}(B)\Longleftrightarrow\operatorname{Skew}(B)$.
\item If $\operatorname{Sym}(B)$, $\operatorname{Skew}(B)$, or $\operatorname{Alt}(B)$, then $\operatorname{rad}_{L}(B)=\operatorname{rad}_{R}(B)$.
\item If $\operatorname{char}(\mathbb{F})=2$, then $\operatorname{Sym}(B)\Longleftrightarrow\operatorname{Skew}(B)$.
\item If $\operatorname{char}(\mathbb{F})=2$, then $\operatorname{Alt}(B)\Rightarrow\operatorname{Sym}(B)$.
\end{enumerate}
\begin{proof}
\begin{enumerate}[label=(\roman*)]
\setlength{\itemsep}{0pt}
\item[$(i)$] Under $\operatorname{Alt}(B)$, $(\forall\,\mathbf{x},\mathbf{y}\in\mathcal{V})\,0=B(\mathbf{x}+\mathbf{y},\mathbf{x}+\mathbf{y})=B(\mathbf{x},\mathbf{y})+B(\mathbf{y},\mathbf{x})$ by bilinearity, so $\operatorname{Skew}(B)$.
\item[$(ii)$] Under $\operatorname{char}(\mathbb{F})\neq2$, $\operatorname{Sym}(B)$, and $\operatorname{Alt}(B)$, the relation $\operatorname{Skew}(B)$ holds by (i).\\
Thus $(\forall\,\mathbf{x},\mathbf{y}\in\mathcal{V})\,2B(\mathbf{x},\mathbf{y})=0$, and $(\forall\,\mathbf{x},\mathbf{y}\in\mathcal{V})\,B(\mathbf{x},\mathbf{y})=0$ by $2\neq0$; therefore $B=0$.
\item[$(iii)$] ($\Rightarrow$) $\operatorname{Alt}(B)\Rightarrow\operatorname{Skew}(B)$ by (i).
\item[$(iii)$] ($\Leftarrow$) If $\operatorname{char}(\mathbb{F})\neq2$ and $\operatorname{Skew}(B)$, then $(\forall\,\mathbf{x}\in\mathcal{V})\,2B(\mathbf{x},\mathbf{x})=0$. Thus $(\forall\,\mathbf{x}\in\mathcal{V})\,B(\mathbf{x},\mathbf{x})=0$ by $2\neq0$, so $\operatorname{Alt}(B)$.
\item[$(iv)$] Under $\operatorname{Sym}(B)\vee\operatorname{Skew}(B)$, $(\forall\,\mathbf{x},\mathbf{y}\in\mathcal{V})\,[B(\mathbf{x},\mathbf{y})=0\Longleftrightarrow B(\mathbf{y},\mathbf{x})=0]$.\\
The same equivalence holds under $\operatorname{Alt}(B)$ by (i), so $\operatorname{rad}_{L}(B)=\operatorname{rad}_{R}(B)$ by Definition \ref{def:bilinear-left-right-radicals}(i).
\item[$(v)$] ($\Rightarrow$) Under $\operatorname{Sym}(B)$, $(\forall\,\mathbf{x},\mathbf{y}\in\mathcal{V})\,B(\mathbf{x},\mathbf{y})=B(\mathbf{y},\mathbf{x})=-B(\mathbf{y},\mathbf{x})$, so $\operatorname{Skew}(B)$.
\item[$(v)$] ($\Leftarrow$) Under $\operatorname{Skew}(B)$, $(\forall\,\mathbf{x},\mathbf{y}\in\mathcal{V})\,B(\mathbf{x},\mathbf{y})=-B(\mathbf{y},\mathbf{x})=B(\mathbf{y},\mathbf{x})$, so $\operatorname{Sym}(B)$.
\item[$(vi)$] If $\operatorname{char}(\mathbb{F})=2$ and $\operatorname{Alt}(B)$, then $\operatorname{Skew}(B)$ by (i) and $\operatorname{Sym}(B)$ by (v).
\end{enumerate}
\hlt{Structural Commentary}: Diagonal expansion controls form types; argument reversal preserves zeros and identifies radicals.\\
\hlt{Key Insight}: The scalar $2$ is invertible off characteristic $2$ and vanishes in characteristic $2$.
\end{proof}
\end{proposition}

Henceforth, $B:\mathcal{V}\times\mathcal{V}\to\mathbb{F}$ is symmetric, skew-symmetric, or alternating.\\
Its left and right radicals agree by Proposition \ref{prop:form-type-relations-by-characteristic}(iv), so a single radical and a single orthogonality notation suffice.

\begin{definition} \label{def:matrix-rank-and-discriminant-of-form} \hlt{Matrix, Rank, and Discriminant of a Form}\\
Let $(\mathcal{V},B)$ be a metric vector space with ordered basis $\beta=(\mathbf{v}_1,\ldots,\mathbf{v}_n)$, where $n\geq1$.
\begin{enumerate}[label=(\roman*)]
\setlength{\itemsep}{0pt}
\item The \hlt{matrix of the form} in $\beta$ is $M_{\beta}\defeq[\innerprod{\mathbf{v}_i,\mathbf{v}_j}]_{i,j}$.
\item The \hlt{rank} of the form is $\operatorname{rank}(M_{\beta})$.
\item The \hlt{discriminant set} is $\Delta(\mathcal{V})\defeq\{r^2\det(M_{\beta}):r\in\mathbb{F}^{\times}\}$.
\item If $\det(M_{\beta})\neq0$, then $\Delta(\mathcal{V})=\det(M_\beta)(\mathbb{F}^{\times})^2$. This set represents the square class of $\det(M_\beta)$ in $\mathbb{F}^{\times}/(\mathbb{F}^{\times})^2$.
\end{enumerate}
\end{definition}

\begin{proposition} \label{prop:metric-form-matrix-recovery-and-congruence} \hlt{Matrices and Invariants under Congruence}\\
Let $n=\dim(\mathcal{V})\geq1$ and let $\beta,\gamma$ be ordered bases of $\mathcal{V}$. Put $C\defeq C_{\beta\leftarrow\gamma}$.
\begin{enumerate}[label=(\roman*)]
\setlength{\itemsep}{0pt}
\item $\innerprod{\mathbf{x},\mathbf{y}}=[\mathbf{x}]_{\beta}^{t}M_{\beta}[\mathbf{y}]_{\beta}$.
\item $M_{\gamma}=C^tM_{\beta}C$.
\item $\{M_{\delta}:\delta\text{ is an ordered basis of }\mathcal{V}\}=\{P^tM_{\beta}P:P\in\operatorname{Mat}(n,n,\mathbb{F})^{\times}\}$.
\item $\operatorname{rank}(M_{\gamma})=\operatorname{rank}(M_{\beta})$.
\item $\Delta(\mathcal{V})$ is independent of the ordered basis used in its definition.
\end{enumerate}
\begin{proof}
\begin{enumerate}[label=(\roman*)]
\setlength{\itemsep}{0pt}
\item[$(i)$] By Proposition \ref{prop:bilinear-coordinate-evaluation}, $\innerprod{\mathbf{x},\mathbf{y}}=[\mathbf{x}]_{\beta}^{t}M_{\beta}[\mathbf{y}]_{\beta}$.
\item[$(ii)$] By Proposition \ref{prop:square-bilinear-form-change-of-basis-congruence}(i), $M_{\gamma}=C^tM_{\beta}C$.
\item[$(iii)$] For $\subseteq$, let $\delta$ be an ordered basis of $\mathcal{V}$. Then $M_{\delta}=C_{\beta\leftarrow\delta}^tM_{\beta}C_{\beta\leftarrow\delta}$ by (ii), and $C_{\beta\leftarrow\delta}\in\operatorname{Mat}(n,n,\mathbb{F})^{\times}$ by Theorem \ref{thm:change-of-basis-matrices}(i).\\
For $\supseteq$, fix $P\in\operatorname{Mat}(n,n,\mathbb{F})^{\times}$. Choose an ordered basis $\delta$ satisfying $P^{-1}=C_{\delta\leftarrow\beta}$ by Theorem \ref{thm:change-of-basis-matrices}(ii).\\
Then $P=C_{\beta\leftarrow\delta}$ by Theorem \ref{thm:change-of-basis-matrices}(i), so $P^tM_{\beta}P=M_{\delta}$ by (ii).
\item[$(iv)$] The matrices $C$ and $C^t$ are invertible by Theorem \ref{thm:change-of-basis-matrices}(i) and Corollary \ref{cor:transpose-invertible-iff}, respectively.\\
Thus $\operatorname{rank}(M_{\gamma})=\operatorname{rank}(M_{\beta})$ by (ii) and Theorem \ref{thm:full-rank-products-preserve-rank}(ii).
\item[$(v)$] Put $D_{\eta}\defeq\{r^2\det(M_{\eta}):r\in\mathbb{F}^{\times}\}$ $(\eta\in\{\beta,\gamma\})$ and $c\defeq\det(C)$. Then $c\in\mathbb{F}^{\times}$ by invertibility of $C$.\\
The identity $\det(M_{\gamma})=\det(C^t)\det(M_{\beta})\det(C)=c^2\det(M_{\beta})$ holds by (ii) and the determinant product and transpose laws.\\
Therefore $D_{\gamma}=D_{\beta}$ by the bijection $r\mapsto rc$ of $\mathbb{F}^{\times}$.
\end{enumerate}
\hlt{Structural Commentary}: Congruence exhausts basis matrices; invertible factors preserve rank and square class.\\
\hlt{Key Insight}: Both arguments transform by the same basis-change matrix.
\end{proof}
\end{proposition}
